\section{Lattices in Dataflow Analysis}\label{sec:lattices}

A dataflow analysis assigns to every program point a value from a lattice $L$
whose partial order $\sqsubseteq$ means ``carries less information than''.
The join $\sqcup$ combines facts that arrive along different control-flow edges.

\begin{definition}\label{def:height}
The \emph{height} of $L$ is the length of the longest strictly increasing chain
$x_0 \sqsubset x_1 \sqsubset \dots \sqsubset x_k$ in $L$.
\end{definition}

If every transfer function is monotone and $L$ has finite height, the iterative
algorithm of~\cite{kildall1973} terminates after at most
$\text{height}(L) \times |\text{nodes}|$ updates, regardless of the order in which
nodes are visited.

\begin{itemize}
  \item For constant propagation, $L$ maps each variable to $\bot$, a constant, or
        $\top$, so its height is $2$ per variable.
  \item For interval analysis the height is infinite, and a \textbf{widening}
        operator is needed to force termination.\footnote{Widening trades
        precision for speed; a later \emph{narrowing} pass can recover some of
        the lost precision.}
\end{itemize}

The fixed point found this way is the \emph{least} one above the initial
assignment, which is the most precise answer the lattice can express; see
Definition~\ref{def:height} and Section~\ref{sec:universal} for the notation.
